% Round 16. Local actual activation of published numerical admissibility.
% The published abelian height theorem is not claimed as a new result.
% Mathematical source independently PASS at root, higher_trace and global_height.
% Reviewed mathematical SHA: bf03e6ba96ac9ff1519dd9283b16febd7fb5b3defd5cb3566b9cc80db24e9a6e.
% This activation changes only these provenance comments.
\section{Local abelian admissibility and actual Chow-motivic transfer}
\label{r16ab:section}

The numerical input in this section is published.  K\"unnemann's
2001 paper already constructs the full classical height on abelian
varieties over number fields, and Round 7 recorded that global
admissibility theorem.  Here we isolate its genuinely local theorem,
valid over an arbitrary complete discretely valued field with finite
residue field, and combine it with the frozen actual numerical detector.
The resulting statement concerns ordinary actual Green forms on their
specified exact support complements.  Its Chow-motivic transfer uses
the new correspondence stability proved in this round.  No novelty
claim is made for the underlying abelian numerical theorem or classical
abelian height.

\subsection{The precise published local input}

Fix a complete discrete valuation ring $R$ with finite residue field
$k$ and fraction field $F$ of characteristic zero.  For a smooth
projective $k$-variety $D$, write
\[
 A^u(D)=\operatorname{im}\bigl(\operatorname{CH}^u(D)_{\mathbf Q}
          \longrightarrow H^{2u}(D_{\bar k},\mathbf Q_\ell(u))\bigr),
 \qquad \ell\ne\operatorname{char}k.
 \tag{AB16.1}\label{r16ab:cycle-image}
\]
These are rational cycle-image groups, not numerical groups and not
the full cohomology groups.  If $\mathcal P/R$ is regular, projective
and strictly semistable of relative dimension $d$, with reduced
special fiber $Y=\bigcup_iD_i$, let $Y^{(1)}$ and $Y^{(2)}$ denote
the disjoint unions of components and their pairwise intersections.
Put
\[
 \begin{split}
 A^u(Y)&=\ker\bigl(\rho:A^u(Y^{(1)})\to A^u(Y^{(2)})\bigr),\\
 A_s(Y)&=\operatorname{coker}\bigl(
       \gamma:A^{d-s-1}(Y^{(2)})\to A^{d-s}(Y^{(1)})\bigr),
 \end{split}
 \tag{AB16.2}\label{r16ab:strata-groups}
\]
where $\rho$ and $\gamma$ are the signed restriction and Gysin maps.
There is a natural surjection
$\operatorname{CH}_s(Y)_{\mathbf Q}\twoheadrightarrow A_s(Y)$.
The component restrictions of a Chow class on $\mathcal P$ lie in
$A^u(Y)$.  Write $i:Y\hookrightarrow\mathcal P$.

\begin{proposition}[Published local abelian input, with its original scope]
\label{r16ab:published}
Let $A/F$ be an abelian variety.  After a finite extension $F'/F$
there is a regular projective strictly semistable model
$\mathcal P/R'$ of $A_{F'}$ for which, in every degree, the natural map
\[
                 A_s(Y)\longrightarrow
              H^{\mathrm{et}}_{2s}(Y_{\bar k'},\mathbf Q_\ell(-s))
 \tag{AB16.3}\label{r16ab:homology-injection}
\]
is injective and the complex
\[
 A_{d+1-u}(Y)\xrightarrow{i^*i_*}A^u(Y)
                 \xrightarrow{\ \cap[Y]\ }A_{d-u}(Y)
 \tag{AB16.4}\label{r16ab:local-exact}
\]
is exact at its middle term.
\end{proposition}

\begin{proof}[Literature provenance]
This is K. K\"unnemann, \emph{Height pairings for algebraic cycles on
abelian varieties}, Ann. Sci. \`Ecole Norm. Sup. (4) 34 (2001),
503--523, Section 5.5, Theorem 5.6 and Theorem 6.6.
Semi-abelian reduction, a finite unramified splitting extension and a
finite ramified extension provide the model in Sections 5.5--5.6.
Section 6.2 explicitly works over a complete DVR with finite residue
field; Theorem 6.6 proves \eqref{r16ab:homology-injection} and
\eqref{r16ab:local-exact} for these local models.  Section 5.8's
erratum explicitly retains the original DVR assertion when the base
is complete.  Thus no descent of $A$ to a number field is needed.

The definitions in Sections 2.3--2.6 are exactly cycle images and
their cokernels as in \eqref{r16ab:cycle-image}--\eqref{r16ab:strata-groups}.
The reduction to a toric degeneration in Section 6 proves the assertion
without assuming numerical equivalence equals homological equivalence
for arbitrary finite-field abelian varieties.
\end{proof}

\begin{lemma}[The local exact complex supplies a numerical model witness]
\label{r16ab:numerical-witness}
Let $\mathcal P/R$ be a regular projective strictly semistable model
of a smooth projective generic fiber $V/F$, of dimension $d$.
Assume \eqref{r16ab:homology-injection} and
\eqref{r16ab:local-exact} in the relevant degrees.
For every geometrically homologically trivial
$z\in\operatorname{CH}^r(V)_{\mathbf Q}$, $1\le r\le d$,
there is an extension $\Theta_0\in\operatorname{CH}^r(\mathcal P)_{\mathbf Q}$
of $z$ whose refined restriction to every $D_i$ is geometrically
homologically trivial, and therefore numerically trivial.
Its pure operational vector-bundle representative is vertically
numerically zero in the precise sense of Round 12 \texttt{ND2}.
\end{lemma}

\begin{proof}
Choose any Chow extension $\Theta$ of $z$, for instance the sum of
the cycle closures.  Let $a\in A^r(Y)$ be its tuple of component
restrictions.  The image $a\cap[Y]\in A_{d-r}(Y)$ is the image of
the specialization of $z$.  Geometric cycle classes commute with
specialization, so its image in geometric $\ell$-adic homology is zero.
The injection \eqref{r16ab:homology-injection} makes
$a\cap[Y]=0$ in $A_{d-r}(Y)$.  Exactness gives
$b\in A_{d+1-r}(Y)$ with $i^*i_*b=a$.

Choose $\beta\in\operatorname{CH}_{d+1-r}(Y)_{\mathbf Q}$ mapping
to $b$, using the stated surjection, and put
\[
                       \Theta_0=\Theta-i_*\beta.
 \tag{AB16.5}\label{r16ab:corrected-witness}
\]
Its generic restriction is still $z$.  Its tuple of component
cycle classes is $a-i^*i_*b=0$.  As $A^r(Y)$ is a subspace of
the direct sum of the component cycle images, each component class
itself is zero.  This proves the asserted homological vanishing,
not just a zero class after pushing from $Y$ to $\mathcal P$.

On the regular model, use the pure rational Chern representative
from Round 12 \texttt{alteration/numerical\_witness\_descent.tex},
equation \texttt{ND1}: a virtual vector bundle $\xi_{\Theta_0}$
with $\operatorname{ch}(\xi_{\Theta_0})\cap[\mathcal P]=\Theta_0$
and all operational character degrees other than $r$ equal to zero.
Every integral vertical $r$-cycle lies in some $D_i$.  Its pairing
with $\operatorname{ch}_r(\xi_{\Theta_0})$ is the numerical pairing
of the refined restriction on that $D_i$ with this cycle, hence zero.
This uses only homological triviality implying numerical triviality.

The first part is also the local correction argument appearing in
K\"unnemann's proof of Proposition 3.3, printed pp.511--512; the
global arithmetic gluing in that proposition is not needed here.
\end{proof}

\begin{theorem}[Full local actual strong G for abelian varieties]
\label{r16ab:local-G}
Let $K/\mathbf Q_p$ be finite and let $A/K$ be an abelian variety.
For every $r$ and every geometrically homologically trivial
$z\in\operatorname{CH}^r(A)_{\mathbf Q}$,
\[
                         \mathfrak o_{A,K}^r(z)=0.
 \tag{AB16.6}\label{r16ab:abelian-obstruction-zero}
\]
In particular, for $1\le r\le\dim A$, each specified rational
cycle $Z$ in class $z$ has a zero-curvature ordinary actual Green
form on precisely $A^{\mathrm{an}}\setminus|Z|^{\mathrm{an}}$,
against every complementary compact actual test over the original $K$.

More generally, if $A$ is defined over a finite extension $L/K$,
the same assertion holds on $A$, viewed as a smooth integral
$K$-variety, for geometric homological triviality on every component.
\end{theorem}

\begin{proof}
Apply Proposition~\ref{r16ab:published} after a finite extension
$F/K$ and Lemma~\ref{r16ab:numerical-witness} to the pullback of $z$.
This supplies a pure numerical model bundle witness on $A_F$.
The map $A_F\to A$ is genuinely finite dominant of complete
native degree $[F:K]$.  The frozen Round 12 finite numerical-model
cover theorem \texttt{r12a:finite-numerical-cover} applies in every
degree and proves \eqref{r16ab:abelian-obstruction-zero}.

For the last assertion first choose $F/L$ with this model.  Regard
it over $K^\circ$.  Degrees on its vertical supports are multiplied
by the residue-field degree, and the scheme-theoretic special-fiber
multiplicities, including ramification, are retained.  Zero numerical
pairings remain zero.  The same frozen theorem now applies to the
finite $K$-morphism $A_F\to A$, so its conclusion is directly in
the original $K$-actual theory.  It is not obtained by relabeling an
$L$-test identity.  Degrees zero and outside the Chow range are
vacuous for the geometric homological kernel.

Finally choose the frozen curved pair for the original specified
cycle $Z$ and subtract a global actual primitive of its curvature.
This is the established sufficient direction to strong G and gives
the exact final support and all actual tests.  No support on an
alteration is pushed to the target.
\end{proof}

\subsection{Transfer through actual Chow correspondences}

\begin{theorem}[A supplied Chow-motive factorization through abelian varieties]
\label{r16ab:motivic}
Let $X/K$ be smooth projective geometrically integral.  Suppose that
after a finite extension $L/K$ there are abelian varieties $A_j/L$
and rational Chow correspondences $u_j,v_j$ with opposite degree
shifts such that
\[
                 [\Delta_{X_L}]=\sum_j v_j\circ u_j
           \quad\text{in }\operatorname{CH}^{\dim X}(X_L\times_LX_L)_{\mathbf Q}.
 \tag{AB16.7}\label{r16ab:chow-factorization}
\]
Then the full geometric homological Chow group of $X$ has zero
actual obstruction in every degree.  Consequently strong G holds
for all its specified geometrically homologically trivial rational
cycles over $K$, with their exact supports and full actual tests.

It is enough, in a fixed degree $r$, to supply
$\sum_jv_{j*}u_{j*}=\operatorname{id}$ on
$\operatorname{CH}^r(X_L)_{\mathbf Q,\mathrm{geom\,hom}}$.
In particular the all-degree hypothesis holds when the rational Chow
motive of $X_L$ is a direct summand of a finite direct sum of Tate
twists of motives $h(A_j)$.
\end{theorem}

\begin{proof}
For a geometric homological class $z$, each $u_{j*}z_L$ is again
geometrically homologically trivial.  By
Theorem~\ref{r16ab:local-G}, its obstruction is zero in the native
$K$-theory of $A_j$.  The new correspondence-kernel theorem
\texttt{r16cov:correspondences} sends these classes to zero
obstruction classes $v_{j*}u_{j*}z_L$ on $X_L$, still over $K$.

Here a correspondence defined on a product over $L$ is used in the
$K$-theory by the open-and-closed immersion
\[
                X_L\times_LA_j\hookrightarrow X_L\times_KA_j.
\]
This immersion selects the diagonal factor of the finite etale
algebra $L\otimes_KL$.  Algebraic correspondence composition on
these factors is precisely composition over $L$, and the diagonal
of $X_L$ lies in the same factor.  Thus
\eqref{r16ab:chow-factorization} is an actual Chow correspondence
identity after this interpretation, with no field-degree divisor
inserted into it.  Summing gives
$\mathfrak o_{X_L,K}(z_L)=0$.

Full finite actual Bott--Chern injectivity for $X_L\to X$ now gives
$\mathfrak o_{X,K}(z)=0$, using the complete degree $[L:K]$.
Apply the exact-support sufficient direction as before.
The assertion for a single Chow degree uses the same equality on
that group's elements.  A motivic inclusion and retraction give
\eqref{r16ab:chow-factorization} by their defining Chow identity.
Neither a merely cohomological decomposition nor an unproved lifting
of numerical projectors is used.
\end{proof}

\begin{corollary}[Arbitrary finite products of curves]
\label{r16ab:curve-products}
Let $C_1,\ldots,C_m/K$ be smooth projective geometrically integral
curves and put $X=\prod_iC_i$.  In every codimension, every
geometrically homologically trivial rational Chow class on $X$ has
zero actual obstruction.  Hence every specified cycle in that full
group satisfies strong G, on its exact support complement and
against all original-field actual tests.  The assertion is not
restricted to exterior products of cycles or to external-product tests.
\end{corollary}

\begin{proof}
After a finite extension choose a rational point on each curve.
The rational Chow-motive identities are
\[
 h(C_i)=\mathbf 1\oplus h^1(C_i)\oplus\mathbb L,
       \qquad h^1(C_i)\simeq h^1(J(C_i)),
 \tag{AB16.8}\label{r16ab:curve-motive}
\]
and $h^1(J(C_i))$ is a direct summand of $h(J(C_i))$.
These are identities for rational equivalence: A. J. Scholl,
\emph{Classical motives}, Proc. Sympos. Pure Math. 55, Part 1
(1994), 163--187, Section 3.2, Proposition 4.5 and
Theorems 5.2--5.3 give the curve projectors, the Picard-motive
comparison and the abelian Chow projectors, respectively.
They also cover the genus-zero case, where the middle summand is zero.

Tensor the split inclusions and retractions in
\eqref{r16ab:curve-motive} and expand the finite direct sum.
Each term is a summand of a Tate twist of the motive of a product
of Jacobians.  That product is an abelian variety; the unit terms
are allowed as the zero-dimensional abelian variety.  This supplies
the exact Chow factorization required in
Theorem~\ref{r16ab:motivic}.  Its native finite descent removes the
auxiliary choice of points and field.
\end{proof}

\paragraph{Height consequence and precise novelty boundary.}
For varieties over a number field satisfying the supplied Chow
factorization, the same algebraic identities hold after completion
and finite local extension.  Theorem~\ref{r16ab:motivic} therefore
puts their full geometric homological Chow groups in the proved
global height domain of \texttt{r16h:global-height}.  This supplies
the actual analytic height construction on that full domain.
The classical existence for abelian varieties and its admissible
Chow-motivic consequences were already known and recorded in Round 7;
this section makes no priority claim and does not identify the new
analytic metric with every previously chosen model metric without
the separate comparison theorem.  General varieties have not been
shown to possess the Chow factorization above.

\paragraph{Primary sources inspected in this round.}
K\"unnemann's original PDF, including its local hypotheses and the
complete-DVR erratum, is
\url{https://www.numdam.org/article/ASENS_2001_4_34_4_503_0.pdf}.
Scholl's author-hosted manuscript is
\url{https://www.dpmms.cam.ac.uk/~ajs1005/preprints/classical_motives.pdf}.
The abelian Chow--K\"unneth theorem originates in C. Deninger and
J. Murre, \emph{Motivic decomposition of abelian schemes and the
Fourier transform}, J. reine angew. Math. 422 (1991), 201--219,
DOI \texttt{10.1515/crll.1991.422.201}; the cited Scholl source
proves the particular comparison and projectors used here.
